\chapter{Limitations, costs and a roadmap}\label{ch:roadmap}

% ------------------------------------------------------------------------------------------------
\section*{What and why}

\Cref{ch:design,ch:module,ch:examples} presented the proposal from the inside. The first gave the
extension protocol added to \micro, the second the module \texttt{(aux kanren lambda)} built on it, and
the third the book's programs running beside their λProlog originals. This chapter looks at the proposal
from the outside, as a reviewer would before a merge. It asks five questions.
\begin{enumerate}
  \item \emph{How do we know it works?} What the test suites check, how a case is written, and what a
    green run does and does not establish (\cref{sec:roadmap-testing}).
  \item \emph{Where does it differ from λProlog?} Nine documented limitations and deviations, each
    pinned by a test case, most with a small run (\cref{sec:roadmap-limits}).
  \item \emph{What does it cost?} The measured overhead on first-order search, after a \texttt{∀°},
    and on λ workloads, and where that time goes (\cref{sec:roadmap-perf}).
  \item \emph{How does it enter the repository?} A staged merge, file by file, with a review order
    (\cref{sec:roadmap-merge}), and six later work items (\cref{sec:roadmap-later}).
  \item \emph{What must the author decide?} The questions that only the maintainer of \aux{} can
    answer (\cref{sec:roadmap-questions}).
\end{enumerate}

Every Scheme snippet below was run against the prototype: micro with the protocol applied, plus the
module. Its output is pasted verbatim. The prototype will live in \texttt{doc/hol-kanren/prototype/}:
the module, \texttt{micro.diff}, \texttt{show.diff}, the design notes \texttt{DESIGN.md}, the tests (with
the two copies of micro's suites, \texttt{micro-tests/} and \texttt{micro-tests-with-lambda/}), the
examples, the benchmarks, and the scripts \texttt{build.sh} and \texttt{run-tests.sh}. After
the merge it moves to \texttt{src/aux.kanren.lambda.scm} and \texttt{src/test/microkanren-lambda*.scm}.
The snippets import \scm{(aux base)}, \scm{(aux stream)}, \micro, \scm{(aux kanren micro show)} and
\scm{(aux kanren lambda)}.

Page numbers \enquote{p.N} in the prototype's comments are PDF pages (see \cref{sec:intro-conventions}).


% ------------------------------------------------------------------------------------------------
\section{How the proposal is tested}\label{sec:roadmap-testing}

\subsection{Three layers of tests}\label{sec:roadmap-layers}

The tests come in three layers, each with its own job.
\begin{description}
  \item[Acceptance (A1--A9).] The acceptance criteria that the three prototypes shared, these cases say what the module must
    do: the scope of \texttt{∀}, hypothetical clauses, type inference, object-level substitution, $\Llam$
    unification, the differences from Ma and Friedman~\cite{MaFriedman2021}, the larger programs of the
    book, and micro's stores on λ-terms. A9 says that micro alone has not changed.
  \item[Regression (R, R2, R3).] The proposal went through three adversarial review rounds, which sent
    thirteen, eight and twelve reports. Every report was reproduced before it was fixed, and every code fix
    came with a regression case; each case of the third round fails on the code before its fix. Several reports turned out to describe the
    same defect, and two were documentation errors or a cost, not code defects. So the rounds left 12, 6
    and 9 cases.
  \item[Limits.] A separate suite, \texttt{microkanren-lambda-limits.scm}, pins what the module does
    \emph{not} do, or does differently from λProlog. Its opening doc string reads \enquote{Known limits of
    (aux kanren lambda), each pinned so that a change of behaviour is noticed.} When one of these cases
    fails, a limitation has changed, perhaps for the better, and the documentation must change with it.
\end{description}

\begin{table}[h]
\centering\small
\caption{The two suites of the module, group by group.}\label{tab:roadmap-suites}
\begin{tabular}{@{}llrr@{}}
\toprule
group & subject & cases & assertions \\
\midrule
A1 & \texttt{∀} scope: escapes, aliases, levels reaching clauses & 9 & 22 \\
A2 & \texttt{⊃} hypothetical clauses: scope, order, sharing; courses, jars, reverse & 8 & 21 \\
A3 & \texttt{∀} and \texttt{⊃} together: type inference, forwards and backwards & 3 & 8 \\
A4 & object-level substitution: cbn, β, copy, subst, existential generalisation & 4 & 9 \\
A5 & $\Llam$: flex--rigid, flex--flex, pruning, 8.G--8.M, non-patterns, η & 10 & 40 \\
A6 & Ma \& Friedman 2021, case by case & 5 & 13 \\
A7 & de Bruijn, prenex, mappred, palindromes, Fischer's CPS & 5 & 12 \\
A8 & micro's stores on λ-terms: \texttt{≠°}, \texttt{absent°}, types, printing, grouping & 11 & 65 \\
A9 & a first-order case inside the suite & 1 & 2 \\
R  & the first review round & 12 & 44 \\
R2 & the second review round & 6 & 31 \\
R3 & the third review round & 9 & 53 \\
\midrule
   & \texttt{microkanren-lambda.scm} & 83 & 320 \\
   & \texttt{microkanren-lambda-limits.scm} & 6 & 13 \\
\bottomrule
\end{tabular}
\end{table}

\subsection{How a case is written}\label{sec:roadmap-case}

A case is a list of assertions \texttt{(}$\assert$\texttt{= expected actual)} from \scm{(aux unittest)}. Every expected
answer was written by hand, from the book or from the logic, and none was captured from a run. This
matters: one prototype had a generated snapshot as its suite, which would have frozen its bugs as
expected answers. Here is the A1 case that tests escapes through an alias:
\begin{scheme}[commandchars=\\\{\}]
((test/A1/alias _)
 ; the escape through a younger variable is caught whatever the order of the bindings and their kind.
 (\ensuremath{\assert}= '() (°->list #f (fresh° (q) (∀° (x) (fresh° (y) (=° q (list y)) (=° y x))))))
 (\ensuremath{\assert}= '() (°->list #f (fresh° (q) (∀° (x) (fresh° (y) (=° y x) (=° q (list y)))))))
 (\ensuremath{\assert}= '() (°->list #f (fresh° (q) (∀° (x) (fresh° (y z) (=° q (list y)) (=° y z) (=° z x))))))
 (\ensuremath{\assert}= '() (°->list #f (fresh° (q) (∀° (x) (fresh° (y) (=° y q) (=° y x))))))
 (\ensuremath{\assert}= '() (°->list #f (fresh° (q) (∀° (x) (=° (vector q) (vector x)))))))
\end{scheme}
Three habits of the suite show here. First, the \emph{same} property is asserted under several goal
orders. Conjunction must be commutative, and many defects found by the reviews broke exactly that: an
answer lost in one order of the goals and found in another. Second, the runner states what is compared.
\scm{(°->list #f g)} gives micro's answer as rebuilding code, constraints included, while
\scm{°->list/datum} gives plain data comparable with \scm{equal?} against an answer quoted from the book.
An error is compared through its message. Third, comments cite the book's page or the review report that
the case pins.

\begin{tutorial}[What the levels do in the second assertion]
Both assertions encode $\exists q\,\forall x\,\exists y.\ q=(y)\wedge y=x$, which has no proof: it would
need $q=(x)$, and $q$ is quantified outside $x$. \Cref{sec:examples-scope} follows the first order, where
\textbf{Check} lowers $y$ when $q:=(y)$ is made and then rejects $y':=x$. Here $\ell(q)=0$ and
$\ell(x)=\ell(y)=1$ (\cref{def:design-levels}), and \scm{(=° y x)} comes first. $\ell(y)=1=\mathit{top}$,
the highest level reached (\cref{def:module-state}), and a variable of level $\mathit{top}$ sees every
eigenvariable that exists, so $y:=x$ is not checked (a fast path, \cref{sec:module-fast}). Then
\scm{(=° q (list y))} binds $q:=(y)$; the check (\cref{def:module-check}) walks $(y)$ to $(x)$ and finds
$\ell(x)=1>\ell(q)=0$. No answer. The two orders fail at different bindings, for the same reason. The
other three assertions vary the route: a chain of two variables, a var--var binding in the other
direction, a vector instead of a list.
\end{tutorial}

\subsection{micro's own suites, run twice}\label{sec:roadmap-a9}

Acceptance item A9 lives mostly outside the module's suite, which holds only one first-order case under
that name. A9 is micro's six existing suites, 125 cases in all, run twice against the edited micro: the
core suite and those of aggregation, untagged terms, show, arithmetic and dates. The first copy is
unchanged. The second differs only in its import line, which also loads the module, so that the
extension is \emph{registered} while first-order code runs. For the core suite:
\begin{answers}
5c5
<   (aux unittest) (aux base) (aux stream) (aux kanren micro) (aux fds sbral))
---
>   (aux unittest) (aux base) (aux stream) (aux kanren micro) (aux kanren lambda) (aux fds sbral))
\end{answers}
The expectations are not touched, answer order included. A seventh suite, the one of union-find, tests
\texttt{(aux kanren unionfind)}, a separate engine that does not import micro, so it was not run.

\subsection{The run}\label{sec:roadmap-run}

The script \texttt{build.sh} compiles the edited micro, its sibling modules (show, arith, date) and the
module with \texttt{csc -O2 -d1 -s -J} into a private repository, \texttt{repo/} beside the script, that
shadows the installed micro. The script \texttt{run-tests.sh} then runs the two suites of the module,
the six suites of micro twice (from \texttt{micro-tests/} and \texttt{micro-tests-with-lambda/}), and the
curated examples, with the private repository first on the search path:
\begin{Verbatim}
cd doc/hol-kanren/prototype
sh build.sh
sh run-tests.sh
\end{Verbatim}
Both scripts set \texttt{CHICKEN\_REPOSITORY\_PATH} to \texttt{repo/} followed by the installed repository,
and \texttt{build.sh} passes \texttt{-I} for the installed headers; both paths are hard-coded to
\texttt{\textasciitilde/.local} and must be adapted on another machine. Nothing is installed. The output of
\texttt{run-tests.sh}, verbatim:
\begin{answers}
test/microkanren-lambda ((ran 83) (failed 0))
test/microkanren-lambda-limits ((ran 6) (failed 0))
micro-tests/microkanren ((ran 24) (failed 0))
micro-tests/microkanren-aggregation ((ran 6) (failed 0))
micro-tests/microkanren-untagged ((ran 7) (failed 0))
micro-tests/microkanren-show ((ran 15) (failed 0))
micro-tests/microkanren-arith ((ran 52) (failed 0))
micro-tests/microkanren-date ((ran 21) (failed 0))
micro-tests-with-lambda/microkanren ((ran 24) (failed 0))
micro-tests-with-lambda/microkanren-aggregation ((ran 6) (failed 0))
micro-tests-with-lambda/microkanren-untagged ((ran 7) (failed 0))
micro-tests-with-lambda/microkanren-show ((ran 15) (failed 0))
micro-tests-with-lambda/microkanren-arith ((ran 52) (failed 0))
micro-tests-with-lambda/microkanren-date ((ran 21) (failed 0))
examples.scm ok
\end{answers}
That is $83+6=89$ cases of the module, plus micro's $24+6+7+15+52+21=125$ cases run twice, $339$ cases
in all, with 0 failures. The aux repository was never written while the prototype was built: its
\texttt{git status} shows only the untracked files it had at the start, and both diffs pass
\texttt{patch --dry-run} against the current \texttt{src/}.

\begin{pitfall}[What a green run does not say]
The suites are example-based. They show that the module gives the book's answers on the book's examples,
and that every defect found by review stays fixed. They do not prove soundness. The guarantees on pattern
problems come from the theorems the algorithm follows: Miller's termination and most-general-unifier
results~\cite{Miller1991}, the soundness and completeness of universe tagging~\cite{NadathurJayaramanKwon1995},
and its levels formulation~\cite{NadathurLinnell2005}. The tests check that the code follows those
algorithms on the cases tried. And a passing limits suite is no good news in itself: it means the
limitations are still there.
\end{pitfall}

% ------------------------------------------------------------------------------------------------
\section{Limitations and deviations from λProlog}\label{sec:roadmap-limits}

λProlog is a typed language with its own implementation of the hohh search. \texttt{(aux kanren lambda)}
is an untyped library over a miniKanren with interleaving search and constraint stores that λProlog
does not have. Some departures follow from that: they are the price of the host. Others are choices that
could be revisited. \Cref{tab:roadmap-limits} lists all nine, with the case that pins each. The
subsections take them in order.

\begin{longtable}{@{}r>{\raggedright\arraybackslash}p{0.52\textwidth}>{\raggedright\arraybackslash}p{0.36\textwidth}@{}}
\caption{The limitations, each pinned by a case.}\label{tab:roadmap-limits}\\
\toprule
& limitation & pinned by \\
\midrule\endfirsthead
\toprule
& limitation & pinned by \\
\midrule\endhead
1 & only dynamic relations take hypotheses; no predicate quantification & \texttt{A2/static-relations} \\
2 & untyped: η only at abstractions, \texttt{@} distinct from lists, β may diverge & \texttt{limits/untyped} \\
3 & non-patterns suspended, never enumerated: floundering reported, not detected & \texttt{limits/no-huet} \\
4 & answer order is micro's interleaving & \texttt{A7/prenex}, \texttt{A3/backwards} \\
5 & \texttt{≠°}, \texttt{absent°} and types are extensions, with a chosen reading of \texttt{∀} & \texttt{A8/diseq-extensional}, \texttt{R2/extensional-absent-order} \\
6 & host escapes: levels police logic variables, not Scheme values & \texttt{limits/host-escape} \\
7 & binder names are hints; some projected constraints keep a working variable & \texttt{A6/alpha}, \texttt{limits/absent-projection} \\
8 & costs: no indexing of E and N; the scope check after any \texttt{∀°} & \texttt{bench/} (\cref{sec:roadmap-perf}) \\
9 & the answer variable must be created outside every \texttt{∀°} & \texttt{R2/answer-inside-forall} \\
\bottomrule
\end{longtable}

\subsection{Only dynamic relations take hypotheses}\label{sec:roadmap-lim-dynamic}

In λProlog every predicate is dynamic: \lp{D => G} may add clauses for any predicate, including one
that has static clauses in the program. Here only a relation defined with \scm{define-dynamic-relation}
consults the context Γ, and assuming a clause for any other relation is an error. A \texttt{⊃°} that
silently ignored such a clause would be worse, and all three prototypes did that.
\begin{scheme}
(define-relation (r° x) (=° x 1))
(define-dynamic-relation (d° x) (=° x 1))
(error-message (λ () (°->list #f (fresh° (q) (⊃° ((r° 2)) (r° q))))))
(°->list/datum (fresh° (q) (⊃° ((d° 2)) (d° q))))
\end{scheme}
\begin{answers}
"not a dynamic relation, see define-dynamic-relation"
(2 1)
\end{answers}
The dynamic relation answers 2 before 1, because hypotheses come first, most recent first, then the
static body (\book{78}, \cref{sec:hh-rules}).

Two consequences follow, both because clauses are Scheme closures and not formulas. A program cannot
build a clause as a term and then assume it. And there is no quantification over predicates, so the
book's \lp{reverse}, which hides its auxiliary \lp{rev} under \lp{pi rev\ ...} \book{125--126}
(\cref{def:hoprog-hohhplus}), cannot be transcribed as it stands. The suite runs instead the \fohh{} \lp{reverse} of
Figure~3.4 \book{87}, which assumes the clauses of its auxiliary predicate but does not hide its name: here
\scm{rev°} is a dynamic relation declared at top level, with no static clauses (\texttt{A2/reverse}). Roadmap item (e) addresses
predicate quantification (\cref{sec:roadmap-later-pred}).

\subsection{Untyped terms}\label{sec:roadmap-lim-untyped}

λProlog's terms are simply typed, and its unifier works on βη-long normal forms. The module has no
types. It keeps terms η-short and η-expands one side only when the other is an abstraction, which is
Miller's treatment of the untyped case~\cite{Miller1991}. Three facts follow.
\begin{scheme}
(°->list #f (fresh° (q) (=° (@ 'g 'a) '(g a))))
(μkanren-datum (ƛ (x) (@ 'g x)))
\end{scheme}
\begin{answers}
()
g
\end{answers}
\begin{itemize}
  \item λ-tree application \scm{(@ 'g 'a)} is not the list \scm{'(g a)}: the first is a term of the
    meta-language, the second Scheme data of the object language. Keeping them apart is what keeps
    conjunction commutative. One prototype conflated them and lost that property.
  \item $\abs{x}{(g\,x)}$ is built η-short as $g$.
  \item Full β through \scm{@} on instantiated heads is not guaranteed to terminate on untyped
    meta-terms. With types, β-reduction would terminate.
\end{itemize}
There is also no check that a solution exists in a possibly empty type (\cref{sec:lambda-solutions}).
Roadmap item (f) addresses types (\cref{sec:roadmap-later-types}).

\subsection{Non-patterns are suspended, never enumerated}\label{sec:roadmap-lim-huet}

Outside $\Llam$ the module does what Teyjus does: it suspends the equation in the store E, retries it after every
binding in case it has become a pattern, and reports it in the answer if it is still suspended (\cref{sec:hounif-pragmatics}).
It does not run Huet's procedure (\cref{sec:hounif-huet}). Two runs show what this means:
\begin{scheme}
(°->list #f (fresh° (q F) (=° q F) (=° (@ F 'a) 'b) (=° (@ F 'a) 'c)))
(°->list #f (fresh° (q F) (=° q F) (=° (@ F 'a) '(g a a))))
\end{scheme}
\begin{answers}
((λ (α) (assert (≐ (@ α 'a) 'b)) (assert (≐ (@ α 'a) 'c)) α))
((λ (α) (assert (≐ (@ α 'a) (cons 'g (cons 'a (cons 'a '()))))) α))
\end{answers}
The second problem is the book's Example 8.A (\cref{ex:hounif-nomgu}), $F\,a=g\,a\,a$, written with list
data for the right-hand side. It has four unifiers and no most general one \book{213}. The module does
not choose: it answers that $F$ is any α satisfying the residual. That answer is correct but conditional.

The first problem shows the cost of not choosing. $F\,a=b$ and $F\,a=c$ have no common solution, since
$b\ne c$. Huet's procedure would find that out, but the module only suspends both equations and returns
an answer whose residuals are jointly unsatisfiable. This is \emph{floundering reported but not detected}:
the answer says what was left undecided, and it is up to the reader, or a later binding, to find out that
it is empty. Roadmap item (d) would add an optional Huet search over E (\cref{sec:roadmap-later-huet}).

\subsection{Answer order is micro's interleaving}\label{sec:roadmap-lim-order}

λProlog searches depth-first, left to right; micro interleaves the branches of a disjunction. Where the
book lists answers, the two orders often agree: augmentation order \book{78}, the existential
generalisations \book{203}, the palindromes \book{142}, reverse, the course database (whose extra,
duplicate proofs come in the interleaving's order), and the CPS term
\book{257}. They differ on \lp{prenex}, whose five answers \book{196--197} come out as the book's first,
third, second, fourth and fifth:
\begin{answers}
((all (ƛ (x) (all (ƛ (x1) (and (q x x) (q x x1))))))
 (all (ƛ (x) (all (ƛ (x1) (and (q x1 x1) (q x x1))))))
 (all (ƛ (x) (all (ƛ (x1) (all (ƛ (x2) (and (q x x) (q x1 x2))))))))
 (all (ƛ (x) (all (ƛ (x1) (all (ƛ (x2) (and (q x1 x1) (q x x2))))))))
 (all (ƛ (x) (all (ƛ (x1) (all (ƛ (x2) (and (q x2 x2) (q x x1)))))))))
\end{answers}
The difference also helps. Where depth-first search loops, interleaving finds answers that λProlog misses.
Asked for terms of type $i\to i$, \lp{typeof} run backwards loops under depth-first search. Here it
enumerates terms:
\begin{answers}
((abs (ƛ (x) x))
 (abs (ƛ (x) (app (abs (ƛ (x1) x1)) x)))
 (abs (ƛ (x) (app (abs (ƛ (x1) x)) x))))
\end{answers}
The suite compares answer \emph{sets}, through a sorted copy, wherever interleaving rather than logic
decides the order. Where it compares a sequence against the book, the book's order is the claim being tested.

\subsection{Disequality, absence and types under \texttt{∀}}\label{sec:roadmap-lim-readings}

λProlog has no disequality, no \emph{absento} and no type constraints. micro has all three, so the
module had to decide what they mean when an eigenvariable is involved. For \scm{≠°} there are two
defensible readings of $\forall x.\ q\ne x$.
\begin{definition}[two readings of an eigenvariable under $\ne$]\label{def:roadmap-readings}
Under the \emph{generic} reading, an eigenvariable is the new constant $c\notin\Sig$ introduced by
\rulename{generic} (\cref{fig:rules-right}), distinct from every term the older variables can denote. Under the \emph{extensional} reading, an eigenvariable stands for an
arbitrary term of its universe~\cite{Moiseenko2019,JinEtAl2021}.
\end{definition}
The generic reading is the default: it matches λProlog, where \lp{pi x\ not (q = x)} succeeds, and it
behaves like αKanren's noms~\cite{ByrdFriedman2007}. The extensional reading is chosen by the parameter
\scm{μkanren-λ-reading}:
\begin{scheme}
(°->list #f (fresh° (q) (∀° (x) (≠° q x))))
(parameterize ((μkanren-λ-reading 'extensional)) (°->list #f (fresh° (q) (∀° (x) (≠° q x)))))
(parameterize ((μkanren-λ-reading 'extensional))
  (list (°->list #f (fresh° (q) (∀° (x) (fresh° (y) (≠° y x) (=° y 'a))) (=° q 'ok)))
        (°->list #f (fresh° (q) (∀° (x) (fresh° (y) (≠° y x) (=° y `(f ,x)))) (=° q 'ok)))))
(°->list #f (fresh° (q) (∀° (x) (fresh° (y) (≠° y x) (=° y 'a))) (=° q 'ok)))
\end{scheme}
\begin{answers}
((λ (α) α))
()
(() ((λ () 'ok)))
((λ () 'ok))
\end{answers}
Read them as statements about all $x$.
\begin{itemize}
  \item Generic, $\forall x.\,q\ne x$: $q$ is older than $x$ and can never denote it, so the
    disequality holds for good and vanishes from the answer.
  \item Extensional, $\forall x.\,q\ne x$: $x$ ranges over every term, $q$'s value included, so the
    disequality cannot hold for all $x$ and the goal fails.
  \item Extensional, $\forall x\,\exists y.\ y\ne x\wedge y=a$: the instance $x=a$ refutes it.
    With $y=(f\;x)$ it holds, since no term $x$ equals the list $(f\;x)$ (occurs check).
  \item Generic, $\forall x\,\exists y.\ y\ne x\wedge y=a$: $a$ is not the new name $x$, so it holds.
\end{itemize}
Both readings were implemented so that the order of a binding and a disequality does not change the
answer. \texttt{A8/diseq-extensional} pins both orders of $\forall x.\,F\,x\ne x$ and $F:=\abs{z}{c}$, and
\texttt{R2/extensional-absent-order} pins four goal orders of $\forall x\,\exists Y.\ Y\ne a\wedge
\mathit{absent}(a,Y)\wedge Y=x$, which fails in every order.

The reading has edges, and the proposal documents them. \scm{absent°} and the types always use the generic
reading. A disequality keeps the reading in force when it is posted, with one exception. A first-order
disequality sits in micro's store D, which records no reading. If it later \emph{migrates} to the
module's store N, because one of its sides came to hold a λ-term or an eigenvariable, it takes the reading
in force at the migration. The two readings then differ only if a program changes
\scm{μkanren-λ-reading} between the two goals. Under the extensional reading, once a \texttt{∀°} has run,
a first-order disequality goes to N at once, with its reading.

\subsection{Host escapes}\label{sec:roadmap-lim-host}

Levels police logic variables. A Scheme value that host code carries out of a scope is not a logic
variable, and the levels know nothing about the route it took. The limits suite pins both halves of the
story:
\begin{scheme}
(define stash #f)
(°->list #f (fresh° (q) (∀° (x) (project° ((v x)) (begin (set! stash v) ✓°)))
                        (λ (s) ((=° q stash) s))))
(°->list #f (fresh° (q) (∀° (x) (project° ((v x)) (begin (set! stash v) ✓°)))
                        (∀° (y) (fresh° (z) (=° z stash) (=° q 'leaked)))))
\end{scheme}
\begin{answers}
()
((λ () 'leaked))
\end{answers}
\begin{tutorial}[Why the first escape is caught and the second is not]
In both goals the first \scm{∀°} raises the depth to 1 and creates $x$ with $\ell(x)=1$. \scm{project°}
stores the eigenvariable record in a Scheme variable, and the scope exits, so the depth is back to 0.
\begin{enumerate}
  \item First goal: $q$ was created at depth 0. Binding $q:=x$ is checked, and $\ell(x)=1>\ell(q)=0$ is
    a scope failure. The smuggled name cannot reach an \emph{older} variable.
  \item Second goal: a second \scm{∀°} raises the depth to 1 again, and $z$ is created inside it with
    $\ell(z)=1$. Binding $z:=x$ passes the check, since $\ell(x)=1\le\ell(z)$. A level is a depth, not the
    identity of a scope, so two sibling scopes at the same depth look alike.
\end{enumerate}
\end{tutorial}
Inside the logic no such program can be written: a variable created in a later scope never sees the
earlier eigenvariable, because it is out of its lexical scope. The escape needs \scm{project°} and
\scm{set!}. The point of the first case is that even this route cannot violate the quantifier prefix of
the older variables. That guarantee has a price, discussed in \cref{sec:roadmap-perf-forall}.

\subsection{Binder names are hints; projected constraints}\label{sec:roadmap-lim-hints}

A λ-record stores an interned hint, the name the user wrote, for printing. Terms built alike are
therefore \scm{equal?}, and they group in \scm{groupby°}. α-variants written with different names are
not \scm{equal?}, although they are α-equivalent, which \scm{≐} and \scm{μkanren-α=?} decide:
\begin{scheme}
(list (equal? (ƛ (x) `(f ,x)) (ƛ (x) `(f ,x)))
      (equal? (ƛ (x) `(f ,x)) (ƛ (y) `(f ,y)))
      (≐ (ƛ (x) `(f ,x)) (ƛ (y) `(f ,y))))
\end{scheme}
\begin{answers}
(#t #f #t)
\end{answers}
A hint that \scm{equal?} ignored could not be stored in the record. The proposal chose to print answers
with the user's binder names, and gave up \scm{equal?} modulo α for that (\cref{sec:roadmap-questions}).

The second half of this item is about answers. An \emph{absento} on a flexible term whose head is in
the answer is kept whole, and an argument out of the answer prints as a working variable:
\begin{scheme}
(°->list #f (fresh° (q) (fresh° (Y) (absent° 'c (@ q Y)))))
(°->list #f (fresh° (q) (fresh° (F) (absent° 'c (@ F q)))))
\end{scheme}
\begin{answers}
((λ (α) (assert (absent? 'c (@ α _1))) α))
((λ (α) α))
\end{answers}
The first answer says: there is some \texttt{\_1} such that $c$ does not occur in $\alpha\,\_1$. It cannot
be weakened to an absento on $\alpha$ alone: $\alpha:=\abs{z}{z\,c}$ contains $c$, yet with
$\_1:=\abs{v}{b}$ the application reduces to $b$. In the second answer the head $F$ is out of the answer
and may be a constant function, so the constraint is dropped. A disequality with such an argument,
\scm{(deny (equal? β (@ α _3)))}, and an absento whose datum may still reduce to an atom of the answer,
\scm{(absent? (@ α 'b) 'a)}, are likewise kept in the answer (\texttt{R3/diseq-unreachable-argument},
\texttt{R3/absento-datum-reduces}).

\subsection{Costs}\label{sec:roadmap-lim-costs}

The engine pays for three things that λProlog implementations avoid. Every binding retries all the
suspended equations of E and the disequalities of N when they are not empty, with no index by variable
(the wake-up loop of \cref{sec:module-wake}).
The scope check walks bound terms below the top level from the first \texttt{∀°} on, for the rest of the
search, even after that \texttt{∀°} has exited. And deeply nested binders scale quadratically.
\Cref{sec:roadmap-perf} gives the figures.

\subsection{The answer variable outside every \texttt{∀°}}\label{sec:roadmap-lim-answer}

\scm{°->list} reifies variable 0, the first one created. If it is created inside a \scm{∀°} and its
value holds that scope's eigenvariable, no answer can be written without the eigenvariable. λProlog
reports no binding for a variable quantified inside \lp{pi x\ }. The module raises an error instead of
printing an eigenvariable free:
\begin{scheme}
(error-message (λ () (°->list #f (∀° (x) (fresh° (q) (=° q x))))))
(°->list #f (∀° (x) (fresh° (q) (=° q 'a))))
(°->list #f (fresh° (q) (∀° (x) (=° q 'a))))
\end{scheme}
\begin{answers}
"the answer holds the eigenvariable of a ∀° it was created in: create the variable it reifies outside every ∀°"
((λ () 'a))
((λ () 'a))
\end{answers}
An answer that does not hold the eigenvariable is reported as usual, wherever its variable was created.
The idiomatic form is the third one: create the query variable first.

% ------------------------------------------------------------------------------------------------
\section{Performance}\label{sec:roadmap-perf}

All figures are in milliseconds, compiled with \texttt{csc -O2 -d1}, the best of three runs per round, with
two rounds. Three binaries are compared. \emph{before} is the installed micro, unedited. \emph{after} is
the edited micro with no extension loaded. \emph{ext} is the edited micro with \texttt{(aux kanren
lambda)} loaded, so its hooks are registered. The first-order benchmark is the one the engineering review
used to compare the three prototypes.

\subsection{The protocol on first-order search}\label{sec:roadmap-perf-protocol}

\begin{table}[h]
\centering\small
\caption{The cost of the protocol on first-order search.}\label{tab:roadmap-perf-protocol}
\begin{tabular}{@{}lrrr@{}}
\toprule
benchmark (ms, round 1 / round 2) & before & after & ext \\
\midrule
\texttt{append°} backwards, 2000 answers, full & 6068 / 5923 & 6121 / 5959 & 6151 / 5929 \\
the same, search only & 672 / 662 & 690 / 669 & 691 / 668 \\
open \texttt{append°}, 1000 answers, search & 1311 / 1314 & 1356 / 1326 & 1327 / 1328 \\
perm6 with 15 pairwise \texttt{≠°}, 720 answers & 176 / 178 & 184 / 178 & 192 / 199 \\
\bottomrule
\end{tabular}
\end{table}

Read \cref{tab:roadmap-perf-protocol} column by column. From \emph{before} to \emph{after}, the protocol
alone costs at most 1\% end to end on \texttt{append°} (5923 to 5959 ms in round 2), within the 2.5\%
round-to-round noise of the unedited binary. That is what one \scm{pair?} test per call site costs when
no extension is registered. All call sites are off micro's hottest path: the variable clauses and the
pair--pair clause of unification are untouched.

From \emph{after} to \emph{ext}, \texttt{append°} does not move, but perm6 does: 192 / 199 ms, 8--12\% more
than the unedited micro's 176 / 178 ms (4--12\% more than \emph{after}). perm6 keeps
15 disequalities in D throughout, so after every \scm{=°} the \texttt{post} hook must decide whether a
new binding brought a λ-term or an eigenvariable into one of D's prefixes (\cref{def:micro-prefix}),
which would then have to migrate to N (\cref{sec:roadmap-lim-readings}). That look costs a few microseconds per binding while D is non-empty.
The d2 prototype, on which the proposal is based, paid 45--50\% here.

\subsection{With D or A non-empty}\label{sec:roadmap-perf-DA}

The third review round found that this look was too expensive in another case: one unrelated constraint
before a long first-order search. Recall micro's stores (\cref{sec:micro-stores}): D holds the
disequalities, A the type constraints such as \scm{symbol°}, and T the absentos. With D or A non-empty,
the proposal walked every new binding through the substitution, and \texttt{append°} backwards binds
long ground tails at every step. That walk is the factor of 1.65 to 1.76 in the middle column below. The last row is a control: its constraint goes to T,
which this walk never concerned.
\begin{table}[h]
\centering\small
\caption{One unrelated constraint before the search (\texttt{t} is not used by \texttt{append°}).}\label{tab:roadmap-perf-DA}
\begin{tabular}{@{}lrrr@{}}
\toprule
\texttt{append°} backwards, 2000 answers, search (ms) & after & ext, before the fix & ext, now \\
\midrule
plain & 673 / 669 & 665 & 657 / 664 \\
+ \texttt{symbol° t} & 674 / 675 & 1097 & 687 / 695 \\
+ \texttt{≠° t 1} & 655 / 657 & 1172 & 670 / 675 \\
+ \texttt{absent° t} & 661 / 664 & --- & 656 / 666 \\
\bottomrule
\end{tabular}
\end{table}
The fix bounds the look: the hook inspects the new bindings for at most 64 nodes, and past that it walks
D and A themselves. The order matters. Walking D and A at every \scm{=°}, the first fix tried, cost perm6
60\% (288 / 302 ms), because perm6's 15 prefixes are walked each time. With the bounded look first,
perm6 stays at 193 / 194 ms, against 181 / 179 for \emph{after}.

\subsection{After a \texttt{∀°}}\label{sec:roadmap-perf-forall}

The second review round measured first-order search that runs after some \texttt{∀°} has run and exited.
In \cref{tab:roadmap-perf-forall}, \texttt{pre} is \scm{(∀° (z) ✓°)}, which binds nothing and exits at
once, and \texttt{typeof°} is one type inference of $\abs{x}{x}$ run before the search.
\begin{table}[h]
\centering\small
\caption{First-order search after a \texttt{∀°} has run. A second run gave 680 / 1125 / 1154 and 1354 / 3055 ms; a third, after the third review's fixes, 663 / 1101 / 1090 and 1310 / 2918 ms.}\label{tab:roadmap-perf-forall}
\begin{tabular}{@{}lrrr@{}}
\toprule
benchmark (ms) & no \texttt{∀°} before & \texttt{pre} before & \texttt{typeof°} before \\
\midrule
\texttt{append°} backwards, 2000 answers, search & 663 & 1104 (+66\%) & 1092 (+65\%) \\
open \texttt{append°}, 1000 answers, search & 1296 & 2940 (+127\%) & --- \\
\bottomrule
\end{tabular}
\end{table}

\begin{tutorial}[Where the time goes after a \texttt{∀°}]
Take one step of \texttt{append°} backwards on the list $L$ of 1999 numbers. It binds a variable to a
pair whose tail is a ground suffix of $L$.
\begin{enumerate}
  \item \emph{No \texttt{∀°} has run.} The level table is empty, no scope check runs, and the only walk of
    the bound term is micro's occurs check.
  \item \emph{\texttt{pre} has run.} The depth went to 1 and back to 0, so the highest level reached is
    $\mathit{top}=1$. The search's variables are at level 0, below $\mathit{top}$. Each binding of such a variable gets the
    scope check, whose allocation-free \scm{μkanren-plain?} walks the bound term through the
    substitution, beside micro's occurs check. On \texttt{append°} the two walk the same ground tail of the
    list, so each binding walks that tail twice.
\end{enumerate}
Why not skip the check when the depth is back to 0 and nothing of a higher level is reachable? For logic
variables that would be safe: an exited scope reaches a depth-0 binding only through constraints, which
are checked when they wake. But an eigenvariable smuggled out by host code would then reach an older
variable unchecked, which is the first case of \cref{sec:roadmap-lim-host}. The cheap fix gives up that
guarantee, so it was not made. The proper fix is one walk that does both checks (\cref{sec:roadmap-later-walk}).
\end{tutorial}

\subsection{λ workloads}\label{sec:roadmap-perf-lambda}

\Cref{tab:roadmap-perf-lambda} compares the final module with the three prototypes. The prototypes'
columns are the engineering review's measurements on the same programs.
\begin{table}[h]
\centering\small
\caption{λ workloads. d1 is the layered prototype, d2 the hooks prototype the proposal grew from, d3
the native one.}\label{tab:roadmap-perf-lambda}
\begin{tabular}{@{}lrrrr@{}}
\toprule
benchmark (ms) & final & d2 & d3 & d1 \\
\midrule
\texttt{typeof} of the Church numeral 400 & 254 & 441 & 105 & 385 \\
\texttt{copy} of the Church numeral 400 & 256 & 442 & 116 & 314 \\
\texttt{typeof} of 150 nested binders & 1306 & --- & 1409 & --- \\
\texttt{copy} of 150 nested binders & 1745 & --- & 1779 & --- \\
\texttt{typeof} backwards, 30 (term, type) pairs & 10 & 11 & --- & --- \\
one pattern equation $F\,x_1\dots x_{50}=g\,x_{50}\dots x_1$ under 50 \texttt{∀} & 2 & 6 & --- & --- \\
\bottomrule
\end{tabular}
\end{table}

The final module is 1.7 times faster than d2, from which it grew, and 2.4 times slower than d3 on the
Church numeral. Its history on \texttt{typeof} of Church 400 shows where the time went:
\begin{center}\small
\begin{tabular}{@{}lr@{}}
\toprule
step & ms \\
\midrule
d2 ported, with instantiation that shifts only open terms & 463 \\
allocation-free \scm{μkanren-plain?} for the scope check & 330 \\
no check for bindings of variables at the top level & 311 \\
the \texttt{atom?} hook: micro's occurs check stops rebuilding eigenvariables & 254 \\
\bottomrule
\end{tabular}
\end{center}
The last step is worth explaining. micro's occurs check rebuilds a foreign record as the list of its
slots to look inside it. An eigenvariable holds no logic variable, so the hook lets micro skip it, and
that alone saved 18\%.

A profile of the remaining time puts it in micro's walks, \scm{μkanren-state-find} and its sbral
lookup. These are the var--var chains that the type variables of \texttt{typeof} form, since micro binds
the older variable to the newer, plus the occurs check of every binding of a subterm. Both are on micro's
hot path, outside the extension's reach. d3's lead comes from its in-core binding, which does the occurs
check and the level check in one walk. That is roadmap item (a). The remaining figures show no gap: on
150 nested binders the final module is 2--8\% faster than d3. The module's cost on nested binders grows
quadratically with the nesting (\cref{sec:roadmap-lim-costs}).

The fixes of the review rounds cost nothing measurable here. After the second round: typeof Church 400
260 ms, copy 259 ms, nested typeof 1292 ms, nested copy 1732 ms. After the third: 253, 256, 1284 and
1729 ms.

% ------------------------------------------------------------------------------------------------
\section{A staged merge into the aux repository}\label{sec:roadmap-merge}

The proposal is three independent pieces of different risk. They should enter the repository in three
stages, each reviewed and tested on its own.

\begin{table}[h]
\centering\small
\caption{Which file goes where.}\label{tab:roadmap-files}
\begin{tabular}{@{}lll@{}}
\toprule
stage & from \texttt{doc/hol-kanren/prototype/} & to \\
\midrule
1 & \texttt{micro.diff} (+235/−41) & \texttt{src/aux.kanren.micro.scm} \\
2 & \texttt{aux.kanren.lambda.scm} (1255 lines) & \texttt{src/aux.kanren.lambda.scm} \\
  & \texttt{test/microkanren-lambda.scm}, \texttt{-limits.scm} & \texttt{src/test/microkanren-lambda*.scm} \\
  & \texttt{test/programs.scm} & \texttt{src/test/programs.scm} \\
  & one \texttt{extension} line & \texttt{src/aux.egg} \\
3 & \texttt{show.diff} (+42/−12) & \texttt{src/aux.kanren.micro.show.scm} \\
\bottomrule
\end{tabular}
\end{table}

\subsection{Stage 1: the protocol}\label{sec:roadmap-stage1}

The diff is self-contained and mentions neither λ-terms, levels nor clauses. It could carry another
store as well, nominal terms or finite domains, so it is worth having even if stage 2 is postponed.
Both diffs were checked with \texttt{patch --dry-run} against the current \texttt{src/}. Applying the first
and, once the egg is rebuilt and installed, running micro's suites (without installing anything, the
private repository of \texttt{build.sh} and \texttt{run-tests.sh}, \cref{sec:roadmap-run}, gives the
same run):
\begin{Verbatim}
cd src
patch aux.kanren.micro.scm < ../doc/hol-kanren/prototype/micro.diff
make test-microkanren test-microkanren-show test-microkanren-arith test-microkanren-date
\end{Verbatim}
These four targets run 112 cases. The aggregation and untagged suites have no target of their own: the
\texttt{test} target runs them with the other four (\texttt{src/Makefile}, lines 46--51), or they run alone
with \texttt{cd test \&\& csi -s microkanren-aggregation.scm}, and likewise for
\texttt{microkanren-untagged.scm}. The expected result is \cref{sec:roadmap-run}'s: 125 cases, unchanged answers. The benchmark's
\emph{after} column (\cref{tab:roadmap-perf-protocol}) is the performance criterion. The diff already
documents each hook in micro's comments; the merge should add to them a short section on the protocol as
a whole.

What to review, in this order:
\begin{enumerate}
  \item \emph{The state.} A seventh field \texttt{X}, the extensions' slot, which micro never reads. It
    is copied at micro's two positional constructions (\scm{μkanren-state-update} and \scm{freshª}),
    \scm{μkanren-state-with} accepts \scm{X:}, and \scm{make-μkanren-state} keeps its six-argument
    signature, so the positional states of \texttt{test/microkanren.scm} still work.
  \item \emph{The registry and the call sites.} \scm{make-μkanren-extension} and
    \scm{μkanren-extension-register!}; each call site costs one \scm{pair?} test when the registry is
    empty. There are thirteen hooks: \texttt{term?}, \texttt{atom?}, \texttt{unify}, \texttt{occurs},
    \texttt{post}, \texttt{diseq}, \texttt{tag}, \texttt{parts}, \texttt{value}, \texttt{repr},
    \texttt{terms}, \texttt{forms} and \texttt{environment}; \cref{ch:design} gives their contracts.
  \item \emph{The two invariants} that every hook must keep (\cref{def:design-invariants}): a hook that
    changes nothing answers the very state it was given, and a hook that binds returns the extended
    \texttt{assocs}.
  \item \emph{The two exported primitives}: \scm{μkanren-state-forget}, which unbinds a variable so that a
    \texttt{post} hook can turn a binding just made back into a constraint before D, A and T see it, and
    \scm{μkanren-post-unify}, the post-\scm{=°} chain that an extension goal unifying on its own must call
    (\cref{sec:design-forget}).
  \item \emph{One change of micro's own behaviour.} micro decided an atom against an absento's datum with
    \scm{equal?} only, so \scm{(absent° D q) (=° q 'a) (=° D 'a)} succeeded in micro alone. The diff
    decides it by unification and watches it while pending. This is a bug fix of micro, independent of λ,
    and deserves its own look.
\end{enumerate}

\subsection{Stage 2: the module and its suites}\label{sec:roadmap-stage2}

Besides \texttt{srfi-1} and CHICKEN's own modules, the module imports \scm{(aux base)}, \scm{(aux stream)},
\scm{(aux fds sbral)} and \micro, and registers
its extension when it is loaded. The egg gets one line after micro's, in the style of its neighbours:
\begin{Verbatim}
  (extension aux.kanren.micro)
  (extension aux.kanren.lambda)
  (extension aux.kanren.arith)
\end{Verbatim}
Other entries of \texttt{src/aux.egg} state their dependencies with \texttt{component-dependencies}. Writing
\texttt{(component-dependencies aux.kanren.micro)} would make this one explicit too. The suite
includes \texttt{programs.scm} from its own directory, so the file goes beside it in
\texttt{src/test/}. The Makefile would follow its existing pattern: two lines in \texttt{test}, after
micro's suites, and a target of its own.
\begin{Verbatim}
	cd test && ${CSI} -s microkanren-lambda.scm
	cd test && ${CSI} -s microkanren-lambda-limits.scm
\end{Verbatim}
Then micro's six suites should be run once more with the module loaded, as in
\cref{sec:roadmap-a9}. The simplest way is to keep the second copy of the six suites under
\texttt{doc/hol-kanren/prototype/}, since they differ from \texttt{src/test/} only in their import line.

What to review, in this order. First the limits suite, 51 lines that state what is not promised. Then
the scope check and variable elimination, which carry soundness (\cref{ch:module}). Then the discharge
at \texttt{∀°} exit, then the store N. The printing code comes last: a mistake there shows up in answers,
not in the search.

\subsection{Stage 3: the show module}\label{sec:roadmap-stage3}

\texttt{show.diff} teaches \scm{(aux kanren micro show)} to read \scm{ƛ}, \scm{@}, \scm{≐}, \texttt{∃} and
\texttt{∀} in answers, draws λ-terms in the mermaid cells view with a node per binder, and adds
\scm{°->pretty}. It reads answer forms syntactically and does not import the module, so it is independent
of stage 2. It is useful only with it.
\begin{Verbatim}
cd src
patch aux.kanren.micro.show.scm < ../doc/hol-kanren/prototype/show.diff
make test-microkanren-show
\end{Verbatim}

% ------------------------------------------------------------------------------------------------
\section{Later work}\label{sec:roadmap-later}

Six items were left out of the proposal on purpose. Each should be measured separately, against the
figures of \cref{sec:roadmap-perf} and the suites of \cref{sec:roadmap-testing}. For each, the paragraphs
below say how it would fit the protocol and which concept of the book it serves.

\subsection{(a) One walk for the occurs and the scope check}\label{sec:roadmap-later-walk}

\emph{What.} micro's binding site, \scm{μkanren-state-update}, runs the occurs check. The module's
\texttt{post} hook then walks the same term again for the scope check. A combined walk does both, as d3's
in-core \texttt{occurs/fast} does.

\emph{How it fits.} A new hook at the binding site, which returns a \emph{verdict} instead of micro's
boolean: no occurrence and nothing out of scope; an occurs failure; or a term that needs variable
elimination. With no extension registered, micro keeps its own check.

\emph{What it buys.} Most of the cost of first-order search after a \texttt{∀°} (+65\% to +127\%,
\cref{tab:roadmap-perf-forall}) and the remaining factor of 2.4 against d3 on Church 400
(\cref{tab:roadmap-perf-lambda}). It keeps the host-escape guarantee, because the check still runs at depth 0.

\emph{Book concept.} The mixed prefix of a unification problem (\cref{sec:lambda-unifprob}), read as
universe levels (\cref{sec:rules-levels}, \cref{sec:hounif-pragmatics}). Nadathur, Jayaraman and Kwon's binding rule~\cite{NadathurJayaramanKwon1995}
is one condition on the bound term, so it is natural to check it in the same walk as the occurs check.

\subsection{(b) Var--var orientation by level}\label{sec:roadmap-later-varvar}

\emph{What.} When two variables are unified, bind the one of higher level to the one of lower level.
micro binds the older to the newer, so a variable of a deep scope may become the value of an outer one,
and the scope check must then lower it (\cref{def:rules-binding}).

\emph{How it fits.} It needs a hook in micro's var--var clause, the hottest path after pair--pair. That
is why the graft from d3 was declined. \textbf{Check} (\cref{def:module-check}) already lowers the
younger variable, which gives the same most general unifier up to renaming, so this item is about cost only: fewer lowering bindings, and
perhaps shorter var--var chains, which the profile of \cref{sec:roadmap-perf-lambda} places among the
remaining costs. It must be measured on first-order search too, where the hook would cost on every
var--var binding.

\emph{Book concept.} Raising and lowering in a mixed prefix (\cref{sec:lambda-raising}). Binding the
deeper variable to the outer one is always within scope. Binding the outer variable to the deeper one is
allowed only after the deeper one has been lowered, which is the extra work this item saves.

\subsection{(c) Indexing E and N by variable}\label{sec:roadmap-later-index}

\emph{What.} Today every \scm{=°} retries all suspended equations of E and all disequalities of N, oldest
first, in rounds, until no round binds (\cref{sec:roadmap-lim-costs}). An index from each variable to the
constraints that mention it would retry only those the new bindings touch.

\emph{How it fits.} Entirely inside the extension. E and N live in the slot \texttt{X}, which micro never
reads, and the \texttt{post} hook already receives the new bindings as \texttt{assocs}. No change to micro.
The care needed is in what \enquote{touches} means: a binding wakes a constraint on a variable it reaches
through the substitution, not only on the variable it binds, and the variables created by lowering and
pruning must be indexed too.

\emph{Book concept.} Dynamic $\Llam$ (\cref{sec:hounif-pragmatics}): an equation that does not respect the
$\Llam$ condition is deferred \enquote{until such time that substitutions resulting from solving other parts of
the problem have altered its status} \book{226}. An index decides cheaply which deferred equations a
substitution may have altered.

\subsection{(d) Huet pre-unification as an optional search}\label{sec:roadmap-later-huet}

\emph{What.} An opt-in goal that takes the suspended equations of E and searches their solutions with
Huet's imitation and projection steps~\cite{Huet1975}, each choice a branch of the interleaving search.
On Example 8.A it would enumerate the four unifiers of \cref{ex:hounif-nomgu}. On the first problem of
\cref{sec:roadmap-lim-huet} it would find no solution and fail, so floundering would be detected, not
only reported.

\emph{How it fits.} E is the natural agenda: it already holds the non-pattern equations, closed over the
eigenvariables they mention, and woken equations are retried by the existing loop. Huet's branching is
don't-know nondeterminism, which a miniKanren stream represents directly. It should be a goal the user
calls, not the default, because the search may not terminate (\cref{sec:hounif-hard}).

\emph{The obstacle.} Projection needs the target type and arities, and imitation the arity and argument
types of the constant: \enquote{Huet's procedure cannot run without types} (\cref{sec:hounif-huet}).
So this item presupposes (f).

\emph{Book concept.} Huet's procedure and its matching trees \book{214--220}, and the limit of
pattern unification that \cref{sec:hoprog-panacea} discusses.

\subsection{(e) Predicate quantification}\label{sec:roadmap-later-pred}

\emph{What.} Goals such as \lp{pi rev\ (... rev ...) => rev L K nil}, which introduce a hidden predicate
local to one call \book{125--126}.

\emph{How it fits.} Through dynamic relations created at run time. Γ is an alist keyed by relation, so a
\texttt{∀} over a predicate would allocate a new key, a new dynamic relation with no static clauses, and
pass it to the body. Clauses assumed for it are then visible only inside that body, and no program text
outside can name it. This is a sketch: the proposal does not implement it.

\emph{Book concept.} \hohh$^+$ (\cref{def:hoprog-hohhplus}): a clause head may be a variable bound by an
essentially universal quantifier. By the time the clause is assumed, $\forall R$ has replaced the variable
by a new constant, so the operational behaviour does not change. Only the visibility of names does.

\subsection{(f) Simple types}\label{sec:roadmap-later-types}

\emph{What.} Optional simple types on λ-terms and on logic variables.

\emph{How it fits.} Through the interplay of the \texttt{atom?} and \texttt{tag} hooks, which the proposal
names for this item without designing it. One reading: a type is a tag on a term, and a typed variable can
be η-expanded by its type. Unification could then work on η-long forms,
as λProlog does, instead of η-expanding only when one side is an abstraction
(\cref{sec:roadmap-lim-untyped}).

\emph{What it buys.} Type-directed η, β that is guaranteed to terminate, the existence check on
possibly empty types, and with them item (d).

\emph{Book concept.} The simply typed terms of \cref{ch:lambda} and the βη-long normal forms on which
the book's unification algorithms are stated (\cref{sec:lambda-normal}).

% ------------------------------------------------------------------------------------------------
\section{Open questions for the author}\label{sec:roadmap-questions}

The proposal makes a choice wherever one was needed to build it and test it. Some of those choices are
the maintainer's to make, and are listed here as questions. Each says what the proposal does now and
what is at stake.

\begin{enumerate}
  \item \emph{The module's name.} The proposal keeps \texttt{(aux kanren lambda)}. The layer is λProlog's
    logic, and λKanren is the name the literature gives this line of work, Ma and Friedman's paper
    included~\cite{MaFriedman2021}. A name such as \texttt{(aux kanren hohh)} would describe only the
    goals, while the module also gives λ-tree syntax and $\Llam$ unification. The names inside follow
    \aux: goals end in \texttt{°}, their applicative forms in \texttt{ª}, term constructors have no
    suffix (\scm{ƛ}, \scm{@}), and internals carry the \texttt{μkanren-} prefix.
  \item \emph{The extensional reading.} \scm{μkanren-λ-reading} exposes a second meaning of
    \scm{≠°} under \texttt{∀°} (\cref{def:roadmap-readings}). It costs code and one edge case,
    the reading of a migrated disequality (\cref{sec:roadmap-lim-readings}). Keep it public, keep it as
    an internal option, or remove it and keep only λProlog's generic reading?
  \item \emph{Should micro own the protocol?} The proposal puts a generic extension protocol in micro,
    a global registry, and one client. The alternatives were a layered module that leaves micro untouched
    (d1, which leaked eigenvariables through micro's own \scm{=°}) and λ-terms native in micro (d3, a
    931-line diff that special-cases every walker). A per-state hook list instead of a global registry
    would require every initial state to carry the hooks, since \scm{°->list} starts from
    \scm{μkanren-state-empty}. Is a protocol in micro, with its contracts and invariants, something the
    maintainer wants to keep stable?
  \item \emph{The cost after a \texttt{∀°}.} First-order search after any \texttt{∀°} takes 1.7 to 2.3
    times as long (\cref{tab:roadmap-perf-forall}). The cheap fix gives up the host-escape guarantee
    (\cref{sec:roadmap-lim-host}); the proper one is roadmap item (a). Should stage 2 be merged now and
    (a) done later, or should (a) be a condition of the merge? And is it worth this price to catch
    escapes made with \scm{project°} and \scm{set!}?
  \item \emph{Which relations are dynamic.} In λProlog every predicate takes hypotheses; here only those
    made by \scm{define-dynamic-relation}, and \texttt{⊃°} on any other is an error
    (\cref{sec:roadmap-lim-dynamic}). Should \scm{define-relation} itself become dynamic when the module
    is loaded? The proposal does not measure what that would cost.
  \item \emph{α-equivalence and \scm{equal?}.} Hints are interned, so terms built alike are
    \scm{equal?}, but α-variants with different names are not (\cref{sec:roadmap-lim-hints}). Truly
    nameless records would make \scm{equal?} decide α, at the price of answers printed with generated
    names instead of the user's. Which matters more for the intended use?
  \item \emph{An answer inside \texttt{∀°}.} Reifying a variable created inside a \texttt{∀°} that holds its
    eigenvariable is an error, where λProlog reports no binding (\cref{sec:roadmap-lim-answer}). Keep the
    error, or answer with no binding as λProlog does?
  \item \emph{A defect of micro left as it is.} Under a type of the user's that holds of pairs, micro's
    subsumption turns \scm{absent° '(a) q} into \scm{q ≠ (a)}, so
    \scm{(absent° '(a) q) (nonnum° q) (=° q '((a)))} succeeds: the subsumption assumes that typed values
    are atoms. It involves no λ-term, so the proposal did not touch it. Fix it in micro, with the protocol
    or separately?
  \item \emph{Two clause syntaxes.} Hypothetical clauses may be written infix, \scm{(H ← B)}, or prefix,
    \scm{(← H B)}; both are accepted. Keep both, or settle on one?
\end{enumerate}

% ------------------------------------------------------------------------------------------------
\section{Summary}\label{sec:roadmap-summary}

\begin{itemize}
  \item The module has two suites: 83 cases with 320 hand-written assertions (acceptance A1--A9 and the
    regressions R, R2, R3 of three adversarial review rounds), and 6 cases that pin its limits. micro's six
    suites, 125 cases, pass with the edited micro both without and with the module loaded, answer order
    unchanged: 339 cases, 0 failures.
  \item The suites test behaviour on examples, under several goal orders; the guarantees on pattern
    problems come from the theorems of Miller, of Nadathur, Jayaraman and Kwon, and of Nadathur and
    Linnell, which the algorithm follows.
  \item Nine limitations are documented and pinned. Only dynamic relations take hypotheses, and there is
    no predicate quantification. Terms are untyped. Non-patterns are suspended, never enumerated, so
    floundering is reported, not detected. Answer order is micro's interleaving. \scm{≠°},
    \scm{absent°} and types use the generic reading of \texttt{∀}, with an extensional option for
    \scm{≠°}. Levels police logic variables, not Scheme values. Binder names are hints. Constraints are
    retried unindexed, and the scope check runs after any \texttt{∀°}. The answer variable must be
    created outside every \texttt{∀°}.
  \item The protocol costs about 1\% on first-order search, and the loaded extension 8--12\% on a
    search that keeps many disequalities. After a \texttt{∀°}, first-order search takes 1.7--2.3 times as long,
    because the scope check walks the bound term beside the occurs check. On λ workloads the module is
    1.7 times faster than the d2 prototype it grew from and 2.4 times slower than the native d3 on
    Church 400, the rest of the time being in micro's walks.
  \item The merge has three stages. First the protocol, into micro's source. Then the module, its
    suites, \texttt{programs.scm} and one egg line. Last the show extension. Review the protocol's invariants
    and its one change to micro's semantics first.
  \item Six later items: one walk for the occurs and scope checks; var--var orientation by level; an
    index of E and N; an optional Huet search, which needs types; predicate quantification through dynamic
    relations made at run time; simple types.
  \item Nine questions are the author's to answer, among them the module's name, the extensional reading,
    whether micro should own a protocol, and whether the cost after a \texttt{∀°} must be paid down before
    the merge.
\end{itemize}
